\chapter{Opis koncepcji rozwiązania}
\label{cha:koncepcja}

\section{Ogólna idea}

System MultiVerify realizuje klasyczny pipeline biometryczny podzielony na dwa główne etapy: \textbf{rejestrację (enrollment)} oraz \textbf{identyfikację (1:N)}. W obu przypadkach przetwarzane są dwie modalności równolegle; decyzja końcowa zapada po fazie fuzji wyników cząstkowych.

\section{Schemat blokowy}

Na rys.~\ref{fig:pipeline} przedstawiono przepływ danych w systemie.

\begin{figure}[H]
    \centering
    \begin{tikzpicture}[
        node distance=8mm and 12mm,
        block/.style={rectangle, draw, rounded corners, align=center, minimum width=28mm, minimum height=9mm, font=\small},
        arrow/.style={-{Stealth[length=2mm]}, thick}
    ]
        \node[block] (in) {Zdjęcie +\\nagranie};
        \node[block, below=of in] (reg) {Rejestracja\\(\texttt{register.py})};
        \node[block, below left=14mm and 8mm of reg] (face) {Ekstrakcja\\twarzy};
        \node[block, below right=14mm and 8mm of reg] (voice) {Ekstrakcja\\głosu};
        \node[block, below=18mm of face] (embF) {Baza\\embeddingów twarzy};
        \node[block, below=18mm of voice] (embV) {Baza\\embeddingów głosu};
        \node[block, below=22mm of reg] (probe) {Probe\\(\texttt{main.py})};
        \node[block, below=of probe] (cmp) {Porównanie 1:N\\(cosine similarity)};
        \node[block, below=of cmp] (fus) {Fuzja\\(\texttt{FusionEngine})};
        \node[block, below=of fus] (dec) {Decyzja\\dostęp TAK/NIE};

        \draw[arrow] (in) -- (reg);
        \draw[arrow] (reg) -- (face);
        \draw[arrow] (reg) -- (voice);
        \draw[arrow] (face) -- (embF);
        \draw[arrow] (voice) -- (embV);
        \draw[arrow] (embF) -| (cmp);
        \draw[arrow] (embV) -| (cmp);
        \draw[arrow] (probe) -- (cmp);
        \draw[arrow] (cmp) -- (fus);
        \draw[arrow] (fus) -- (dec);
    \end{tikzpicture}
    \caption{Schemat blokowy systemu MultiVerify.\label{fig:pipeline}}
\end{figure}

\section{Etap rejestracji}

Dla każdego użytkownika system:
\begin{enumerate}
    \item wczytuje pliki \texttt{face\_*} i \texttt{voice\_*} z katalogu \texttt{data/users/<user>/},
    \item wyznacza embeddingi twarzy (wektor Facenet) oraz wektory MFCC głosu,
    \item zapisuje wyniki w plikach \texttt{face\_embeddings.npy} oraz \texttt{voice\_embeddings.npy}.
\end{enumerate}

\section{Etap identyfikacji}

Podczas identyfikacji system:
\begin{enumerate}
    \item ekstrahuje cechy z nowego zdjęcia i nagrania (probe),
    \item dla każdego użytkownika w bazie oblicza wynik podobieństwa twarzy i głosu,
    \item wybiera najlepszy wynik w ramach wielu wzorców jednego użytkownika,
    \item łączy wyniki w silniku fuzji,
    \item wskazuje użytkownika z najwyższym wynikiem końcowym i weryfikuje próg akceptacji.
\end{enumerate}

\section{Parametry decyzyjne}

Domyślnie zastosowano:
\begin{itemize}
    \item wagi fuzji: $w_f = 0{,}6$, $w_v = 0{,}4$,
    \item próg akceptacji: $\tau = 0{,}7$,
    \item strategia domyślna: suma ważona (\texttt{weighted}).
\end{itemize}

Parametry te mogą być modyfikowane z linii poleceń programu \texttt{main.py}.
